
\begin{ledgroupsized}[r]{120mm}
\footnotesize
\pstart
\noindent\textbf{Überlieferung:}
\pend
\end{ledgroupsized}
\begin{ledgroupsized}[r]{114mm}
\footnotesize
\pstart \parindent -6mm
\makebox[6mm][l]{\textit{L}}%
Konzept:
LH~XXXV~10, 13~Bl.~1\textendash2. 
Ein Bogen~2\textsuperscript{o};
\href{\Scan{LH_35_10_13_002wz}}%
{Wasserzeichen} in Bl.~2 und 
\href{\Scan{LH_35_10_13_001wz}}%
{Gegenmarke} in Bl.~1;
alle Ränder beschnitten.
Dreieindrittel Seiten.
\pend
\end{ledgroupsized}
\selectlanguage{latin}
\vspace{8mm}
\pstart%
\noindent%
\edtext{\href{\Scan{LH_35_10_13_001r}}%
{\lbrack1~r\textsuperscript{o}\rbrack}}{%
\applabel{FnA:35_10_13_001-002_Bernoullius}
\lemma{\textit{In der oberen linken Ecke nachgetragen:}}%
\Afootnote{%
Addantur quae 
\indexN{\textls{Bernoulli}, Jacob 1655\textendash1705}Bernoullius\lbrack,\rbrack\
\indexN{\textls{Eschinardi}, Francesco 1623\textendash1703}Eschinardus 
\lbrack et\rbrack\textsuperscript{\lbrack a\rbrack} alii
in \cite{01023}\textit{Actis Eruditorum}\textsuperscript{\lbrack b\rbrack} %
\lbrack\textit{Text bricht ab.}\rbrack
\\[1.5ex] %
{\footnotesize 
\textsuperscript{\lbrack a\rbrack} et \textit{erg.~Hrsg.} \quad 
\textsuperscript{\lbrack b\rbrack} in \textit{Actis Eruditorum}: 
Zu Bernoulli siehe 
\indexN{\textls{Bernoulli}, Jacob 1655\textendash1705}\textsc{Jacob Bernoulli},
\cite{02003}\glqq Solutio difficultatis contra propositionem quandam mechanicam Autore I.\,F.\,V. Lucensi propositae\grqq, \cite{01023}\textit{AE} Februar 1686, S.\;96\textendash98.
\indexN{\textls{Eschinardi}, Francesco 1623\textendash1703}F.~Eschinardi reagierte auf Vannis Thesen in seinem \cite{02105}\textit{Cursus physicomathematicus}, tom.~I, Rom 1689, bes. app.~IV, S.\;358\textendash361 (siehe auch die Erläuterung zu \refnote{FnC:AE_1694_112-113_Eschinardi}). 
Die \cite{01023}\textit{Acta Eruditorum} enthalten dennoch keine Publikation Eschinardis zu diesen Themen.
Zu den \glqq alii\grqq\ siehe z.\,B. 
\indexN{\textls{Kochański}, Adam Adamandus 1631\textendash1700}\textsc{A.\,A. Kochański}, %
\cite{01996}\textit{Consideratio speciminis libri de momentis gravium, Authore J.\,F.\,V. Lucensi}, %
\cite{01023}\textit{Acta eruditorum},
Juni 1685, S.\;262\textendash265.}}}
\pend
\nobreak
\begin{edfigure}[0pt][1em]    	%
\hspace{4mm}\includegraphics[scale=1]{%
gesamttex/edit_VIII-4/images/LH_35_10_13_001-002_d1_001r.pdf}
\edcaption{%
\lbrack\textit{Fig.~1}\rbrack}
\end{edfigure}
\pstart
\noindent
\indexN{\textls{Vanni}, Giovanni Francesco 1638\textendash1709}Johannes Franciscus Vannius
Jesuita Romanus
\edtext{librum ante annos complures ediderat 
\cite{02000}\textit{de Momentis Gravium},}{%
\applabel{FnC:LH_35_10_13_001-002-complures}%
\lemma{librum \lbrack...\rbrack\ \textit{Gravium}}%
\Cfootnote{%
\indexN{\textls{Vanni}, Giovanni Francesco 1638\textendash1709}\textsc{G.\,F.~Vanni}, \cite{02000}\textit{Exegeses physico-mathematicae, de momentis gravium...}, Rom 1685. Leibnizens Rezension erscheint als N.\;\ref{41199} in diesem Band.}}
in quo 
\edtext{stabilitam doctrinam\indexS{doctrina stabilita de gravitatione in plano inclinato} evertere conabatur, de gravitatione in plano inclinato.\indexS{planum inclinatum}}{%
\lemma{stabilitam \lbrack...\rbrack\ inclinato}%
\Cfootnote{%
\indexN{\textls{Vanni}, Giovanni Francesco 1638\textendash1709}Vanni kritisiert insbesondere
\indexN{\textls{Galilei} (Galilaeus, Galileus), Galileo 1564\textendash1642}%
\textsc{G.~Galilei}, %
\cite{00050}\textit{Discorsi}, %
Giornata Terza, De motu naturaliter accelerato; \glqq Aggiunta postuma dell'Autore\grqq\ %
nach dem Scholion zum 2.~Theorem (der Zusatz erschien erst in der posthumen Werkausgabe: \cite{01084}\textit{Opere}, Bologna 1656, Bd.~II, S.\;132\lbrack ter\rbrack\textendash134; \cite{00048}\textit{GO} VIII, S.\;216\textendash219 im Apparat) %
und
\indexN{\textls{Torricelli} (Torricellius), Evangelista 1608\textendash1647}%
\textsc{E.~Torricelli}, \textit{De motu gravium}, lib.~I
(\cite{00304}\textit{Opera geometrica}, Florenz 1644, S.\;99;
\cite{00108}\textit{TO}~II, S.\;105).}}
\edtext{Considerabat inter alia corpus \textit{A} incumbens 
\edtext{simul duabus regulis\indexS{corpus duabus regulis incumbens}}{%
\lemma{simul}%
\Bfootnote{%
\textit{(1)}~duobus planis
\textit{(2)}~duabus regulis~\textit{L}}}
rigidis, una \textit{BC}, altera \textit{DE}, ex quibus prior sit verticalis,\indexS{regula verticalis} posterior vero sit inclinata,\indexS{regula inclinata}
et hinc paralogismis 
\edtext{quibusdam intercedentibus impugnabat doctrinam receptam\indexS{doctrina stabilita de gravitatione in plano inclinato} de gravitatione in plano inclinato\indexS{planum inclinatum} novamque}{%
\lemma{quibusdam}%
\Bfootnote{%
\textit{(1)}~adhib
\textit{(2)}~intercedentibus
\textit{(a)}~novam
\textit{(b)}~impugnabat \lbrack...\rbrack\ novamque~\textit{L}}}
plane 
\edtext{Staticen\indexS{Statica nova} procudere}{%
\lemma{Staticen}%
\Bfootnote{%
\textit{(1)}~componere
\textit{(2)}~procudere~\textit{L}}}
volebat.}{%
\lemma{Considerabat \lbrack...\rbrack\ volebat}%
\Cfootnote{%
\indexN{\textls{Vanni}, Giovanni Francesco 1638\textendash1709}Vanni, \cite{02000}\textit{Exegeses}, \textit{passim}.
Siehe auch das als Vorankündigung des Buchs verstandene \cite{01995}\textit{Specimen libri de momentis gravium}, \cite{01023}\textit{AE} November 1684, S.\;511\textendash513, das bereits die wesentlichen Beweise enthielt.
Vgl. dazu Leibnizens anonyme, in den \cite{01023}\textit{Acta Eruditorum} zusammen mit dem \textit{Specimen} abgedruckte Anmerkung (N.\;\ref{55606} in diesem Band).}}
\edtext{Ego memini me olim soluisse difficultatem, sed cum 
\edtext{non sit}{%
\lemma{non}%
\Bfootnote{%
\textit{(1)}~satis 
\textit{(2)}~recordari
\textit{(3)}~sit~\textit{L}}}
ad manum vetus scheda,\indexS{scheda vetus}}{%
\applabel{FnC:35_10_13_001-002-scheda}%
\lemma{Ego \lbrack...\rbrack\ scheda}%
\Cfootnote{Gemeint sind die unveröffentlicht gebliebenen Konzepte \textit{Solutio problematis R.\,P.\;Vannii S.\,J.} (April bis Juni 1690), N.\;\ref{41823} in diesem Band.}}
et nuper considerationem renovaverim, repetere solutionem placet.
\pend
\pstart
\edtext{\textls{Theorema 1.} Si}{%
\lemma{\textls{Theorema 1.}}%
\Bfootnote{%
\textit{(1)}~Si
\textit{(a)}~gravis
\textit{(b)}~gravi
\textit{(c)}~globus 
\textit{(d)}~gravitate
\textit{(2)}~Si~\textit{L}}}
globus\indexS{globus gravis ex materia similari}
\edtext{gravis}{%
\lemma{}%
\Bfootnote{%
gravis
\textit{erg.~L}}}
ex materia similari simul 
\edtext{incumbat duabus regulis\indexS{globus duabus regulis incumbens}}{%
\lemma{incumbat}%
\Bfootnote{%
\textit{(1)}~duobus planis
\textit{(2)}~duabus regulis~\textit{L}}}
in eodem plano inter se et cum centro globi\indexS{centrum globi} positis
\edtext{quarum altera sit verticalis,\indexS{regula verticalis}}{%
\lemma{}%
\Bfootnote{%
quarum altera sit 
\textit{(1)}~perpendicu
\textit{(2)}~verticalis
\textit{erg.~L}}}
aequales erunt pressiones regularum,\indexS{pressio regulae}
\edtext{et alterutra erit ad pondus ipsum ut secans anguli complementi regularum.\indexS{secans anguli complementi regularum}}{%
\lemma{}%
\Bfootnote{%
,~et
\textit{(1)}~utraque
\textit{(2)}~alterutra
erit \lbrack...\rbrack\ regularum
\textit{erg.~L}}}
\pend
\begin{edfigure}[1em][1em]    	%
\includegraphics[scale=1]{%
gesamttex/edit_VIII-4/images/LH_35_10_13_001-002_d2_001r.pdf}
\edcaption[0.25em]{%
\lbrack\textit{Fig.~2}\rbrack}
\end{edfigure}
\pstart
Sit globus gravis similaris\indexS{globus gravis similaris} cujus centrum \textit{A}, regulae
\edtext{rigidae,}{%
\lemma{}%
\Bfootnote{%
rigidae
\textit{erg.~L}}}
\textit{BC}, et \textit{DE} quas simul
\edtext{ambas}{%
\lemma{}%
\Bfootnote{%
ambas
\textit{erg.~L}}}
tangit premitque globus illam in \textit{F}, hanc
\edtext{in \textit{G}. Ponatur regula ita coercita}{%
\lemma{in \textit{G}.}%
\Bfootnote{%
\textit{(1)}~Ponantur 
\textit{(a)}~aute
\textit{(b)}~regulae ita coercitae
\textit{(2)}~Ponatur regula ita coercita~\textit{L}}}
ut non possit recedere nisi in normali ad se, simulque in quantum 
\edtext{recedit, in tantum}{%
\lemma{recedit,}%
\Bfootnote{%
\textit{(1)}~\bars{in \textit{streicht Hrsg.}} tantum attrahere pondera 
\textit{(2)}~normaliter
\textit{(3)}~in tantum~\textit{L}}}
verticaliter attrahere pondera \textit{I}, et \textit{K}, quaeritur horum ponderum quantitas.
Jungantur rectae ex \textit{A} regulis
\edtext{\textit{BC}, \textit{DE}}{%
\lemma{}%
\Bfootnote{%
\textit{BC}, \textit{DE}
\textit{erg.~L}}}
normales,
\edtext{nempe}{%
\lemma{}%
\Bfootnote{%
nempe
\textit{erg.~L}}}
\textit{AF}, et 
\edtext{\textit{AG}, et rectae}{%
\lemma{\textit{AG},}%
\Bfootnote{%
\textit{(1)}~et iisdem parallelae \textit{AL} et \textit{L}
\textit{(2)}~et rectae~\textit{L}}}
iisdem regulis parallelae \textit{AL} et \textit{AH}
quarum illa secet
\textit{DE} 
\edlabel{35_10_13_001-002_6a}%
\edtext{}{%
\xxref%
{35_10_13_001-002_6a}{35_10_13_001-002_6b}%
\lemma{in \textit{L},}%
\Bfootnote{%
\textit{(1)}~prior \textit{BC}
\textit{(2)}~haec
\textit{(a)}~\textit{BC} 
\textit{(b)}~\textit{BC} 
in \textit{H}.
\bars{Compleatur \bars{trapezium \textit{erg.~Hrsg.}} rectangulum \textit{AHLG}. \textit{erg.}}
\textit{(aa)}~Posito
\textit{(bb)}~Jam~\textit{L}}}%
in \textit{L}, haec \textit{BC} in \textit{H}. 
\edtext{Compleatur \lbrack trapezium\rbrack\ rectangulum \textit{AHLG}.}{%
\lemma{Compleatur \lbrack trapezium\rbrack\ rectangulum \textit{AHLG}}%
\Cfootnote{%
Das Viereck \textit{AHLG} ist seiner Definition gemäß kein Rechteck (\glqq rectangulum\grqq\ in \textit{L}), sondern ein rechtwinkliges Trapez, da \textit{AH} und \textit{AL} gleich sind.}}
Jam\edlabel{35_10_13_001-002_6b}
si ambae regulae sufficienter resistant, perinde aget grave, ac si essent
immobiles. Si ergo immobilis esset regula \textit{DE},
grave \textit{A} in recta 
\edtext{\textit{A}\lbrack\textit{H}\rbrack}{%
\lemma{}%
\Bfootnote{%
\textit{AL} 
\textit{L ändert Hrsg.}}}
ei parallela conaretur descendere, 
et gravitatio ejus foret ad totalem seu 
\edtext{liberam\indexS{gravitatio gravis totalis seu libera} ut \textit{A}\lbrack\textit{H}\rbrack\ ad \textit{AL}.}{%
\lemma{liberam ut}%
\Bfootnote{%
\textit{(1)}~\textit{AL} ad
\textit{(2)}~\bars{\textit{AC} \textit{ändert Hrsg.}} ad \textit{AL}.~\textit{L}}}
Hoc ergo 
\edtext{nisu ut \textit{A}\lbrack\textit{H}\rbrack\ agit}{%
\lemma{}%
\Bfootnote{%
nisu ut \bars{\textit{AL} \textit{ändert Hrsg.}} agit~\textit{L}}}
in regulam 
\edtext{\textit{BC},}{%
\lemma{}%
\Bfootnote{%
\textit{BC}
\textit{erg.~L}}}
seu regula haec non impellitur nisi in recta \textit{AF} seu 
\edlabel{35_10_13_001-002_5a}%
\edtext{}{%
\xxref{35_10_13_001-002_5a}{35_10_13_001-002_5b}%
\lemma{\textit{Am Rand, gestrichen:}}%
\Afootnote{%
\newline\nobreak
{\footnotesize 
\nobreak
\centerline{\includegraphics[scale=1]{%
gesamttex/edit_VIII-4/images/LH_35_10_13_001-002_d3_001r.pdf}} 
\centerline{\lbrack\textit{Fig.~3}\rbrack}
\\[0.5em]
$AC^2 = AB^2+BC^2$. 
Sit $AC = 2AB$\lbrack,\rbrack\ fit $4AB^2 = AB^2+BC^2$ 
seu $BC=AB\surd 3$.}}}%
\edtext{nisu perpendiculari\indexS{nisus perpendicularis} ad ipsam, qui sit ad nisum parallelum\indexS{nisus parallelus}}{%
\lemma{nisu}%
\Bfootnote{%
\textit{(1)}~qui sit ad \textit{AL}, ut
\textit{(2)}~parallelo ad \textit{DE},
\textit{(3)}~perpendiculari ad ipsam,
\textit{(a)}~seu qui
\textit{(b)}~qui sit ad
\textit{(aa)}~\textit{AH} ut \textit{AF}
\textit{(bb)}~nisum
\textit{(aaa)}~quo
\textit{(bbb)}~parallelum~\textit{L}}}
ipsi \textit{DE}, ut \textit{AF} ad \textit{AH}. Ergo 
\edtext{erit nisus}{%
\lemma{erit}%
\Bfootnote{%
\textit{(1)}~nisus gravi
\textit{(2)}~nisus~\textit{L}}}
quo propelletur regula \textit{BC}, ad nisum gravis liberum,\indexS{nisus gravis liberus} 
\edtext{u\lbrack t\rbrack}{%
\lemma{}%
\Bfootnote{%
u \textit{L ändert Hrsg.}}}
\textit{AF} ad \textit{AL}.
Similiter \textit{A} descendens in recta parallela ipsi regulae \textit{BC}, id est in recta 
\edtext{vertic\lbrack ali\rbrack}{%
\lemma{}%
\Bfootnote{%
verticulari
\textit{L ändert Hrsg.}}}
\textit{AL}, agit in regulam \textit{DE}, secundum normalem \textit{AG},
eritque nisus quo haec regula impellitur, ut \textit{AG} 
\edtext{ad \textit{AL}. Sed}{%
\lemma{ad \textit{AL}.}%
\Bfootnote{%
\textit{(1)}~Sed \textit{AG} et \textit{AH} sunt
\textit{(2)}~Sed~\textit{L}}}
\textit{AF} et \textit{AG} sunt aequales.
Aequales ergo sunt resistentiae regularum seu aequalia erunt pondera\indexS{pondus} \textit{I} et \textit{K}, 
\edtext{et unumquodque erit ad pondus\indexS{pondus} \textit{A} ut \textit{AG} radius globi, ad \textit{AL},
est autem \textit{AG} ad \textit{AL}, ut radius ad secantem anguli \textit{GAL} qui est complementum 
anguli 
\edtext{\lbrack\textit{HAL}\rbrack,}{%
\lemma{}%
\Bfootnote{%
\textit{GAL} 
\textit{L ändert Hrsg.}}}
quem faciunt regulae inter se.\indexS{secans anguli complementi regularum}%
\edlabel{35_10_13_001-002_5b}}{%
\lemma{}%
\Bfootnote{%
,~et 
\textit{(1)}~utrumque erit ad
\textit{(a)}~\textit{H}
\textit{(b)}~\textit{A}, ut radius sphaerae ad
\textit{(2)}~unumquodque erit ad pondus \textit{A} ut 
\bars{\textit{AG} \textit{erg.}} radius \lbrack...\rbrack\ inter se~\textit{erg.~L}}}
\pend
\pstart
Hinc patet etiam pondera\indexS{pondus}
\edtext{\textit{I}, \textit{K}}{%
\lemma{}%
\Bfootnote{%
\textit{I}, \textit{K}
\textit{erg.~L}}}
duo simul sumta non aequari ponderi \textit{A}, sed esse majora\lbrack,\rbrack\
neque enim $AF + AG$ seu duplex radius globi rectae \textit{AL} aequatur.
Sed si poneremus \textit{A} perpendiculariter descendere tantillum, oportet ut 
\edtext{celeritates ipsorum}{%
\lemma{celeritates}%
\Bfootnote{%
\textit{(1)}~ipsarum
\textit{(2)}~ipsorum~\textit{L}}}
\textit{A} et \textit{I} sint reciproce ut pondera, quia tunc regula 
\edtext{\textit{B}\lbrack\textit{C}\rbrack}{%
\lemma{}%
\Bfootnote{%
\textit{BE} 
\textit{L ändert Hrsg.}}}
non impellitur.
Et si \textit{A} descendat in alia linea recta quacunque media inter \textit{AL} et \textit{AF},
\edtext{oportet factum ex \textit{A} in celeritatem perpendiculariter}{%
\lemma{oportet}%
\Bfootnote{%
\textit{(1)}~descensum ejus
\textit{(2)}~celeritatem
\textit{(3)}~factum ex
\textit{(a)}~massa
\textit{(b)}~\textit{A} in celeritatem
\textit{(aa)}~perpendiculare %
\textit{(bb)}~perpendiculariter~\textit{L}}}
aestimati descensus\indexS{celeritas descensi perpendiculariter aestimati} aequari 
\edtext{factis ex pondere unoquoque}{%
\lemma{factis ex}%
\Bfootnote{%
\textit{(1)}~pondera quoque
\textit{(2)}~pondere unoquoque~\textit{L}}}
duorum reliquorum ducto in celeritates perpendiculares\indexS{celeritas perpendicularis} suas, 
quod pro examine nostri theorematis valere posset.
\edtext{Et ex}{%
\lemma{Et}%
\Bfootnote{%
\textit{(1)}~hoc
\textit{(2)}~ex~\textit{L}}}
hoc ipso posset inveniri;
quia duabus assumtis 
\edtext{lineis res}{%
\lemma{lineis}%
\Bfootnote{%
\bars{ut \textit{erg. u. gestr.}}
\textit{(1)}~res 
\textit{(2)}~determinaturque
\textit{(3)}~res~\textit{L}}}
determinatur, et revera nos duas ex his maxime determinatas assumsimus, nempe extremas,
et duae sufficere debent ad alias omnes determinandas, 
alioqui problema foret impossibile. Quod absurdum.
\pend
\begin{edfigure}%
\includegraphics[scale=1]{%
gesamttex/edit_VIII-4/images/LH_35_10_13_001-002_d4_001r.pdf}
\edcaption{%
\edtext{\lbrack\textit{Fig.~4, gestr.}\rbrack}{%
\lemma{}%
\Afootnote[nonum]{%
\textit{Oberhalb \lbrack\textit{Fig.~4}\rbrack, ebenfalls gestrichen:}\enskip %
{\footnotesize
$\iota \lambda + \kappa \mu = A \beta$.
$\iota \lambda + \iota \mu = A \beta$.\textsuperscript{\lbrack a\rbrack}
$\iota = A \beta: \lambda + \mu$.
$2 \iota = A$. Foret
$\iota = 2$
\\[1.5ex] %
\textsuperscript{\lbrack a\rbrack} $\iota \mu = A \beta$. \textit{(1)}~$\iota = \lambda + \mu$ \textit{(2)}~$\iota =$~\textit{L}}}}}
\end{edfigure}
\pstart
Ex hoc ipso ergo patet methodum nostram esse 
\edlabel{35_10_13_001-002_1a}%
\edtext{}{%
\xxref%
{35_10_13_001-002_1a}{35_10_13_001-002_1b}%
\lemma{probam.}%
\Bfootnote{%
\textit{(1)}~\lbrack/\rbrack\ Theorema II
\textit{(2)}~\lbrack/\rbrack\ Vannii~\textit{L}}}%
probam.
\pend
\pstart
Vannii\indexN{\textls{Vanni}, Giovanni Francesco 1638\textendash1709}\edlabel{35_10_13_001-002_1b}
autem peccatum in eo potissimum constitit, 
quod rem secundum ipsum ad pondera\indexS{pondus} reducendo, oporteret 
\edtext{esse $I+K=A$. Quod certo}{%
\lemma{esse}%
\Bfootnote{%
\textit{(1)}~\textit{K}, et
\textit{(2)}~$I+K=A$.
\textit{(a)}~Quod absurdum est,
\textit{(aa)}~quia descendente
\textit{(bb)}~nisi in eo casu
\textit{(aaa)}~quo casus perpendi
\textit{(bbb)}~quo \bars{is \textit{streicht Hrsg.}}  est
\textit{(ccc)}~nisi is situs esset ut descendente \textit{A} in linea
\textit{(aaaa)}~quadam recta
\textit{(bbbb)}~recta qua maxime proficeret aequaliter ascenderent \textit{I} et \textit{K},
\textit{(aaaaa)}~essetque
\textit{(bbbbb)}~posset esse
\textit{(ccccc)}~essetque ascensus
\textit{(aaaaa-a)}~reciproce corporum 
\textit{(bbbbb-b)}~ponderum duorum ad Ascensum globi, ut globus ad pondera.
Quod qui situs efficere potest, dignum inquisitione foret. Nempe descendente \textit{A} in $\alpha$, 
ponatur maxime proficere positis \textit{I} et \textit{K} aequalibus, 
\textit{(aaaaa-aa)}~seu ponatur centrum gravitatis 
\textit{(bbbbb-bb)}~datis, 
\textit{(ccccc-cc)}~seu ponatur centrum gravitatis maxime descendere, videndum an tunc intervalla
inter \textit{BC} et $\beta\zeta$, vel
\textit{(aaaaa-aaa)}~iter
\textit{(bbbb-bbb)}~inter \textit{DE} et $\delta\varepsilon$ sint aequalia.
Vel potius eligatur is descensus ipsius \textit{A} in linea $A\alpha$, tali, ut fiant 
promoties\lbrack\textit{!}\rbrack\ 
regularum aequales, unde haberi potest, qualia tunc debeant esse pondera aequalia \textit{I} et \textit{K}. Quo facto videatur an data hoc modo pondera \textit{I} et \textit{K} sint
\textit{(aaaaa-aaaa)}~capacia
\textit{(bbbb-bbbb)}~sufficientia ad efficiendum, ut \textit{A} non descendat alio modo. Atque ita definiri poterit, quis situs sit talis, ut hoc obtineatur.
\textit{(b)}~Quod certo~\textit{L}}}
tantum 
\edtext{casu contingere}{%
\lemma{}%
\Bfootnote{%
casu \bars{ex demonstrato \textit{gestr.}} contingere~\textit{L}}}
potest.
Nam quia \textit{K} ad \textit{A} ut \textit{AG} ad \textit{AL}, 
\edtext{ergo si}{%
\lemma{ergo}%
\Bfootnote{%
\textit{(1)}~si
\textit{(2)}~qui
\textit{(3)}~si~\textit{L}}}
$I+K=A$ ideo quia \textit{I} et \textit{K} aequales, 
erit $2K=A$ seu \textit{K} ad \textit{A}, ut 1 ad 2,	%
ergo et \textit{AG} ad \textit{AL}, ut 1 ad 2.
Angulus ergo \textit{GAL} debet esse talis, ut secans sit radii 
\edlabel{35_10_13_001-002_2a}%
\edtext{}{%
\xxref%
{35_10_13_001-002_2a}{35_10_13_001-002_2b}%
\lemma{duplus.}%
\Bfootnote{%
\href{\Scan{LH_35_10_13_001v}}%
{\lbrack1~v\textsuperscript{o}\rbrack}
\textit{(1)}~THEOREMA II.
\textit{(2)}~%
\bars{\textls{Theorema 2.}
\textit{(1)}~Pondera quae retinent
\textit{(2)}~Iisdem positis nisi quod neutra regularum sit verticalis,
\bars{quemcumque faciant angulum \textit{erg.}}
erunt pressiones regularum ut secantes angulorum quos faciunt ad verticalem.
\textit{streicht Hrsg.}}
\textit{(3)}~Theorema 2.~\textit{L}}}%
duplus.
\pend
\pstart
\edtext{%
\href{\Scan{LH_35_10_13_001v}}%
{\lbrack1~v\textsuperscript{o}\rbrack}%
}{%
\lemma{}%
\Afootnote{%
\textit{Am oberen Seitenrand nachgetragen, gestrichen:} %
{\footnotesize 
Rad.
globi\textsuperscript{\lbrack a\rbrack}
ad \textit{AN} ut rad.
globi\textsuperscript{\lbrack b\rbrack}
ad ($AL\;\text{sec}\,LAN:\text{rad}$)\textsuperscript{\lbrack c\rbrack} 
seu in ratione composita radii globi ad \textit{AL}, et
radii ad secantem anguli \textit{LAN}. Sed radius globi est ad 
\textit{AL}, ut\textsuperscript{\lbrack d\rbrack}
radius ad secantem complementi anguli regularum.\indexS{secans anguli complementi regularum} Ergo rad. 
globi\textsuperscript{\lbrack e\rbrack}
est ad \textit{AN},
seu\textsuperscript{\lbrack f\rbrack}
nis\lbrack us\rbrack\ ponderis in regulam 
est ad\textsuperscript{\lbrack g\rbrack} 
pondus globi in ratione composita radii ad secantem complementi anguli regularum,\indexS{secans anguli complementi regularum} 
et radii ad secantem anguli quem regula altera facit ad verticalem.\indexS{secans anguli regulae ad verticalem}
\\[1.5ex] %
\textsuperscript{\lbrack a\rbrack} globi \textit{erg.~L}\quad 
\textsuperscript{\lbrack b\rbrack} globi \textit{erg.~L}\quad 
\textsuperscript{\lbrack c\rbrack} $:\text{rad}$) 
\textit{(1)}~\bars{seu ut qu. \textit{streicht Hrsg.}} rad. ad $AL\;\text{sec}\,LAN$ sed \textit{AL} 
\textit{(a)}~est
\textit{(b)}~est ad rad. globi ut rad. globi in rad.
\textit{(2)}~seu~\textit{L}\quad 
\textsuperscript{\lbrack d\rbrack} \textit{AL}, ut
\textit{(1)}~secans
\textit{(a)}~compl
\textit{(b)}~anguli
\textit{(2)}~radius~\textit{L}\quad 
\textsuperscript{\lbrack e\rbrack} globi \bars{\textit{AG} seu nisus \textit{streicht Hrsg.}} est~\textit{L}\quad 
\textsuperscript{\lbrack f\rbrack} seu \bars{nisi \textit{ändert Hrsg.}}
\textit{(1)}~in regu
\textit{(2)}~ponderis~\textit{L}\quad 
\textsuperscript{\lbrack g\rbrack} est ad
\textit{(1)}~nisum absol
\textit{(2)}~pondus globi \bars{ut \textit{streicht Hrsg.}} in~\textit{L}%
}}}
\edlabel{35_10_13_001-002_8}%
Theorema 2.\edlabel{35_10_13_001-002_2b}
Si globus\indexS{globus gravis similaris}
\edtext{gravis similaris}{%
\lemma{gravis}%
\Bfootnote{%
\textit{(1)}~ex ma
\textit{(2)}~similaris~\textit{L}}}
simul incumbat duabus rectis rigidis,\indexS{globus duabus rectis rigidis incumbens} 
sive regulis, in eodem plano inter se, et cum centro globi\indexS{centrum globi} 
\edtext{positis, pressiones regularum\indexS{pressio regulae} (seu}{%
\lemma{positis,}%
\Bfootnote{%
\textit{(1)}~erunt
\textit{(2)}~pressiones regularum
\textit{(a)}~,~ut
\textit{(b)}~(seu~\textit{L}}}
pondera quae retinent regulas) erunt inter se ut secantes angulorum, quos
\edtext{faciunt regulae}{%
\lemma{faciunt}%
\Bfootnote{%
\textit{(1)}~ad
\textit{(2)}~regulae~\textit{L}}}
ad verticalem\lbrack,\rbrack\indexS{secans anguli regulae ad verticalem} et cujusque regulae pressio\indexS{pressio regulae} ad pondus absolutum\indexS{pondus absolutum globi} 
\edtext{globi est}{%
\lemma{globi}%
\Bfootnote{%
\textit{(1)}~ut
\textit{(2)}~est~\textit{L}}}
in ratione 
\edtext{composita ex}{%
\lemma{composita}%
\Bfootnote{%
\textit{(1)}~radii ad se
\textit{(2)}~ex~\textit{L}}}
ratione radii ad secantem complementi anguli regularum inter se,\indexS{secans anguli complementi regularum} 
et radii ad secantem anguli quem altera regula facit ad verticalem.\indexS{secans anguli regulae ad verticalem}
\pend
\begin{edfigure}%
\hspace{-5mm}\includegraphics[scale=1]{%
gesamttex/edit_VIII-4/images/LH_35_10_13_001-002_d5_001v.pdf}
\edcaption{%
\lbrack\textit{Fig.~5}\rbrack}
\end{edfigure}
\pstart
Sit globus\indexS{globus gravis}
\edtext{gravis}{%
\lemma{}%
\Bfootnote{%
gravis
\textit{erg.~L}}}
centro 
\edtext{\textit{A}, tangens}{%
\lemma{\textit{A},}%
\Bfootnote{%
\textit{(1)}~premens reg
\textit{(2)}~tangens~\textit{L}}}
et premens regulam 
\edtext{\textit{B}\lbrack\textit{C}\rbrack}{%
\lemma{}%
\Bfootnote{%
\textit{BE} 
\textit{L ändert Hrsg.}}}
in \textit{F}, et regulam \textit{DE} in \textit{G}, 
per \textit{A} ad \textit{BC} ducatur \textit{AH} parallela ipsi \textit{DE}, et per \textit{A} ad \textit{DE}
ducatur \textit{AL} parallela ipsi 
\edtext{\textit{B}\lbrack\textit{C}\rbrack.}{%
\lemma{}%
\Bfootnote{%
\textit{BE} 
\textit{L ändert Hrsg.}}}
\edtext{Eruntque \textit{AH} et \textit{AL} aequales, nam in triangulis rectangulis \textit{AFH} et \textit{AGL},
sunt aequales \textit{AF} et \textit{AG}, angulique his oppositi \textit{AHF} et \textit{ALG}.}{%
\lemma{}%
\Bfootnote{%
Eruntque \textit{AH} et \textit{AL} aequales, nam 
\textit{(1)}~angulus
\textit{(2)}~in triangulis
\lbrack...\rbrack\ 
et \textit{ALG}
\textit{(a)}~,~et \textit{FH} et \textit{HL}, seu e
\textit{(b)}~.~\textit{erg.~L}}%
\lemma{Eruntque \lbrack...\rbrack\ et \textit{ALG}}%
\Cfootnote{%
Eine durch diese Ergänzung überschriebene Zeichnung eines rechtwinkligen Dreiecks wird nicht wiedergegeben.}}
Sit per \textit{A} verticalis 
\edtext{\textit{AV}, huic ex \textit{H}}{%
\lemma{\textit{AV},}%
\Bfootnote{%
\textit{(1)}~cui
\textit{(2)}~huic
\textit{(a)}~verti
\textit{(b)}~ex \textit{H}~\textit{L}}}
educta \textit{HM} parallela 
\edtext{ipsi \textit{AG}}{%
\lemma{ipsi}%
\Bfootnote{%
\textit{(1)}~\textit{AH}
\textit{(2)}~\textit{AG}~\textit{L}}}
adeoque recta ad \textit{AH}, occurrat in \textit{M}.
Et eidem verticali \textit{AV} 
\edtext{ex \textit{L} educta}{%
\lemma{ex \textit{L}}%
\Bfootnote{%
\textit{(1)}~ad re
\textit{(2)}~educta~\textit{L}}}
\textit{LP} parallela ipsi \textit{AF}, adeoque recta ad \textit{AL}, occurrat 
\edtext{in \textit{P}. Globus}{%
\lemma{in \textit{P}.}%
\Bfootnote{%
\textit{(1)}~Globus agit in
\textit{(a)}~\textit{BC}
\textit{(b)}~regulam \textit{BC}, quasi
\textit{(2)}~Globus~\textit{L}}}
\textit{A} impeditus a regula \textit{DE}, tendit descendere in plano inclinato\indexS{planum inclinatum} seu in recta
\edtext{\textit{AH}}{%
\lemma{}%
\Bfootnote{%
\textit{AH}
\textit{erg.~L}}}
parallela 
\edtext{ipsi \textit{DE}. Et}{%
\lemma{ipsi \textit{DE}.}%
\Bfootnote{%
\textit{(1)}~Et nisum in \textit{AH} esse
\textit{(2)}~Et~\textit{L}}}
gravitationem in \textit{AH} esse ad gravitationem absolutam seu ad pondus\indexS{pondus} \textit{A}, ut \textit{AH} ad \textit{AM}.
\edtext{Sed ex gravitatione in \textit{AH} tantum}{%
\lemma{Sed ex}%
\Bfootnote{%
\textit{(1)}~nisu \textit{A}
\textit{(2)}~sec
\textit{(3)}~gravitationem,
\textit{(4)}~gravitatione in \textit{AH}
\textit{(a)}~ad gravitationem
\textit{(b)}~tantum~\textit{L}}}
pervenit in \textit{BC} pars quae sit 
\edtext{ad totam}{%
\lemma{ad}%
\Bfootnote{%
\textit{(1)}~totum nisum
\textit{(2)}~totam~\textit{L}}}
gravitationem in \textit{AH}, ut \textit{AF} ad \textit{AH}.
Ergo gravitatio qua 
\edtext{premitur regula\indexS{pressio regulae}}{%
\lemma{premitur}%
\Bfootnote{%
\textit{(1)}~\textit{BC}
\textit{(2)}~regula~\textit{L}}}
\textit{BC} a 
\edtext{pondere \textit{A},}{%
\lemma{pondere}%
\Bfootnote{%
\textit{(1)}~\textit{H},
\textit{(2)}~\textit{C},
\textit{(3)}~\textit{A},~\textit{L}}}
est ad pondus\indexS{pondus} \textit{A}, ut \textit{AF} 
\edtext{radius globi}{%
\lemma{}%
\Bfootnote{%
radius globi
\textit{erg.~L}}}
ad 
\edtext{\textit{AM}. Et}{%
\lemma{\textit{AM}.}%
\Bfootnote{%
\textit{(1)}~Porro
\textit{(2)}~Et~\textit{L}}}
eodem modo gravitatio qua premitur regula\indexS{pressio regulae} \textit{DE} a pondere \textit{A}, est ad pondus\indexS{pondus} \textit{A}, 
ut \textit{AG} radius globi ad \textit{AP}, ergo pressio regulae \textit{BC} 
est ad pressionem regulae\indexS{pressio regulae} \textit{DE} ut \textit{AP} ad \textit{AM}. Porro
\textit{AM} est ad \textit{AH} ut secans 
\edlabel{35_10_13_001-002_3a}%
\edtext{}{%
\xxref%
{35_10_13_001-002_3a}{35_10_13_001-002_3b}%
\lemma{anguli \textit{HAM},}%
\Bfootnote{%
\textit{(1)}~qui est angulus regulae \textit{DE} ad verticalem \textit{AV}, 
\textit{(a)}~secans
\textit{(b)}~est ad radium.
\textit{(2)}~est ad \lbrack...\rbrack\ \textit{AV}.
\textit{(a)}~Porro
\textit{(b)}~Et~\textit{L}}}%
anguli \textit{HAM}, est ad radium. Angulus autem \textit{HAM} 
est angulus regulae \textit{DE} cum verticali \textit{AV}.\indexS{secans anguli regulae ad verticalem}
\pend
\vspace{0.5em}
\pstart
\noindent
\lbrack%
\textit{Nachfolgend kleingedruckter Text (erste Fassung von \refpassage{35_10_13_001-002_7a}{35_10_13_001-002_7b}) gestrichen:}\rbrack %
\pend
\nobreak
\vspace{0.5em}
\footnotesize
\pstart 
\noindent
Et\edlabel{35_10_13_001-002_3b}
\textit{AH} est ad \textit{AF}
\edtext{radium globi}{%
\lemma{}%
\Bfootnote{%
radium globi
\textit{erg.~L}}}
ut secans
\edtext{Anguli \textit{HAF} est ad radium. Sed jam angulus \textit{HAF} est angulus complementi anguli \textit{FHA}}{%
\lemma{Anguli \textit{HAF}}%
\Bfootnote{%
\textit{(1)}~(qui est angulus complementi
\textit{(2)}~est \lbrack...\rbrack\
\textit{HAF} est \bars{angulus \textit{erg.}} complementi anguli
\textit{(a)}~\textit{FAH}
\textit{(b)}~\textit{FHA}~\textit{L}}}
seu angulus complementi ejus quem faciunt regulae 
\edtext{inter se,\indexS{angulus complementi regularum} (unde cum \textit{AF} sit aequalis \textit{AG}, 
et angulus \textit{FHA} aequalis angulo \textit{GLA}, patet \textit{AH} et \textit{AL} esse aequales). 
Ex his sequitur radium globi esse ad \textit{AM}, seu gravitationem in regulam \textit{BC}, esse ad pondus}{%
\lemma{inter se,}%
\Bfootnote{%
\textit{(1)}~. Ergo
\textit{(a)}~est
\textit{(aa)}~r
\textit{(bb)}~ra
\textit{(b)}~radius globi est ad \textit{AM}, seu gravitatio in regulam \textit{BC}, est
\textit{(2)}~, (unde cum \textit{AF} sit aequalis \textit{AG},
\textit{(a)}~patet
\textit{(b)}~et angulus \lbrack...\rbrack\ \textit{BC}, esse ad 
\textit{(aa)}~totam gravitat
\textit{(bb)}~pondus~\textit{L}}}
globi \textit{A}, in ratione composita radii ad secantem anguli complementi regularum,\indexS{secans anguli complementi regularum}
et radii ad secantem anguli regulae alterius
\edtext{\textit{DE}}{%
\lemma{}%
\Bfootnote{%
\textit{DE}
\textit{erg.~L}}}
ad verticalem,\indexS{secans anguli regulae ad verticalem} quod erat in theoremate 
\edlabel{35_10_13_001-002_4a}%
\edtext{}{%
\xxref%
{35_10_13_001-002_4a}{35_10_13_001-002_4b}%
\lemma{secundum.}%
\Bfootnote{%
\textit{(1)}~Et gravitationem regulae \textit{BC} esse in gr
\textit{(2)}~Porro quia gravitatio
\textit{(a)}~regulae
\textit{(b)}~in regulam \textit{BC} est ad gravitationem in regulam \textit{DE}, 
ut secantem anguli \textit{DE} 
\textit{(aa)}~ad vert
\textit{(bb)}~\textit{B}
\textit{(cc)}~regulae \textit{BC}
\textit{(aaa)}~ad verticalem
\textit{(bbb)}~cum verticali ad secantem anguli regulae \textit{DE} cum verticali. 
Quod erat demonstrandum. \lbrack/\rbrack\
\textit{(aaaa)}~Rur
\textit{(bbbb)}~Rursus qui
\textit{(cccc)}~Eodem modo
\textit{(aaaaa)}~qui
\textit{(bbbbb)}~erit \textit{AP} ad
\textit{(dddd)}~Porro eodem modo, ut invenimus \textit{AH},
\textit{(eeee)}~Patet et fore
\textit{(ffff)}~Quia ergo eodem modo radius 
\textit{(gggg)}~Porro quia eodem modo radius globi est ad
\textit{(aaaaa)}~\textit{AM}
\textit{(bbbbb)}~\textit{AP} in ratione composita radii ad secantem anguli complementi regularum et radii ad secantem anguli regulae
\textit{(aaaaa-a)}~\textit{BC}
\textit{(bbbbb-b)}~\textit{DE}
\textit{(ccccc-c)}~\textit{BC} ad verticalem, ideo erit
\textit{(aaaaa-aa)}~\textit{AM} ad
\textit{(bbbbb-bb)}~\textit{AP} ad \textit{AM}
\textit{(hhhh)}~Est autem
\textit{(iiii)}~Et ratio \bars{radii \textit{erg.}} globi
\textit{(aaaaa)}~ad
\textit{(bbbbb)}~ad \textit{AP} est ad rationem 
\textit{(aaaaa-a)}~globi ad
\textit{(bbbbb-b)}~radii
\textit{(jjjj)}~Et ratio quam radius globi facit ad \textit{AP}, est ad rationem quam radius globi facit ad \textit{AM}, 
\textit{(aaaaa-a)}~ut \textit{AM} ad \textit{AP},
\textit{(bbbbb-b)}~erit ad\lbrack\textit{!}\rbrack\ \textit{AM} ad \textit{AP}, 
\textit{(ccccc-c)}~ut ratio composita. Et in utraque ratione radii globi ad \textit{AP} et radii globi ad \textit{AM}, communis est ratio radii ad secantem anguli complementi regularum, ergo
\textit{(aaaaa-aa)}~erit \textit{AM} ad \textit{AP} ut ratio 
\textit{(bbbbb-bb)}~ut secantem anguli complementi regularum ad
\textit{(3)}~Porro
\textit{(4)}~Porro quia
\textit{(a)}~\textit{AH}
\textit{(b)}~\textit{A}
\textit{(c)}~\textit{AM} ad
\textit{(d)}~quemad
\textit{(e)}~invenimus \textit{AH} esse ut
\textit{(5)}~Porro~\textit{L}}}%
secundum.
\pend
\pstart
Porro\edlabel{35_10_13_001-002_4b}
quia invenimus esse \textit{AM} ad \textit{AH} ut secans 
\edtext{anguli regulae}{%
\lemma{anguli}%
\Bfootnote{%
\textit{(1)}~\textit{DE}
\textit{(2)}~regulae~\textit{L}}}
\textit{DE} cum 
\edtext{verticali\indexS{secans anguli regulae ad verticalem} est ad}{%
\lemma{verticali}%
\Bfootnote{%
\textit{(1)}~ad
\textit{(2)}~est ad~\textit{L}}}
radium, et similiter \textit{AP} ad \textit{AL} ut secans anguli regulae \textit{BC} cum verticali\indexS{secans anguli regulae ad verticalem}
est ad radium, seu \textit{AL} et \textit{AH} sunt aequales ut supra est ostensum.
Erit \textit{AP} ad 
\edtext{\textit{AM}, id est}{%
\lemma{\textit{AM},}%
\Bfootnote{%
\textit{(1)}~ut secans anguli regulae \textit{BC} cum
\textit{(2)}~id est~\textit{L}}}
(per ostensa) 
\edtext{pressio in regulam}{%
\lemma{pressio}%
\Bfootnote{%
\textit{(1)}~regulae
\textit{(2)}~in regulam~\textit{L}}}
\textit{BC} ad pressionem in regulam \textit{DC}, ut secans anguli regulae \textit{BC}
cum verticali,\indexS{secans anguli regulae ad verticalem} ad secantem anguli regulae \textit{DE} cum verticali.\indexS{secans anguli regulae ad verticalem} Quod erat primum.
\pend
\normalsize %
\vspace{0.5em}
\pstart
\edlabel{35_10_13_001-002_7a}%
\href{\Scan{LH_35_10_13_002r}}%
{\lbrack2~r\textsuperscript{o}\rbrack}\ 
Et eodem modo \textit{AP} est ad \textit{AL} ut secans 
\edtext{Anguli \textit{LAP} est}{%
\lemma{Anguli}%
\Bfootnote{%
\textit{(1)}~\textit{HAM}
\textit{(2)}~\textit{LAP}
\textit{(a)}~seu
\textit{(b)}~\bars{cum verticali, \textit{streicht Hrsg.}}
\textit{(c)}~est~\textit{L}}}
ad radium, estque angulus \textit{LAP}, is quem regula \textit{BC} facit cum verticali.\indexS{secans anguli regulae ad verticalem}
Itaque quia \textit{AL} et \textit{AH} aequales, erit \textit{AP} ad
\edtext{\textit{AM}, adeoque}{%
\lemma{\textit{AM},}%
\Bfootnote{%
\textit{(1)}~ut secans anguli \bars{quem \textit{streicht Hrsg.}} re
\textit{(2)}~adeoque~\textit{L}}}
pressio regulae\indexS{pressio regulae} \textit{BC} ad pressionem regulae\indexS{pressio regulae} \textit{DE}, ut secans anguli 
quem regula \textit{BC} facit cum verticali\indexS{secans anguli regulae ad verticalem} ad secantem anguli quem regula \textit{DE}
facit cum verticali,\indexS{secans anguli regulae ad verticalem} \textls{quod in Theoremate erat primum}. Porro habuimus 
\edtext{rationem \textit{AM}}{%
\lemma{rationem}%
\Bfootnote{%
\textit{(1)}~\textit{AH} ad \textit{AM}
\textit{(2)}~\textit{AM}~\textit{L}}}
ad \textit{AH}. Jam \textit{AH} est ad \textit{AF} radium globi, ut secans anguli \textit{HAF} est 
ad radium. Jam angulus \textit{HAF} est angulus complementi anguli \textit{FHA},
seu angulus complementi ejus anguli quem faciunt regulae inter se.\indexS{angulus complementi regularum}
Itaque \textit{AH} est ad \textit{AF} ut secans complementi anguli regularum\indexS{secans complementi anguli regularum} est ad radium.
Componendo 
\edtext{ergo rationes}{%
\lemma{ergo}%
\Bfootnote{%
\textit{(1)}~rationem
\textit{(2)}~rationes~\textit{L}}}
\textit{AF} ad \textit{AH}, et \textit{AH} ad \textit{AM}, jam inventas;
erit \textit{AF} ad \textit{AM}, seu pressio regulae 
\edtext{\textit{BC} ad pondus \textit{A}, in}{%
\lemma{}%
\Bfootnote{%
\textit{BC} \bars{erit \textit{streicht Hrsg.}} ad pondus \textit{A},
\textit{(1)}~ut
\textit{(2)}~in~\textit{L}}}
ratione composita radii ad secantem complementi anguli regularum,\indexS{secans complementi anguli regularum}
et radii ad secantem anguli regulae alterius 
\edtext{\textit{DE}, cum verticali.\indexS{secans anguli regulae ad verticalem}}{%
\lemma{\textit{DE},}%
\Bfootnote{%
\textit{(1)}~ad verticalem
\textit{(2)}~cum verticali.~\textit{L}}}
\textls{Quod in Theoremate erat secundum.}%
\edlabel{35_10_13_001-002_7b}%
\pend
\begin{edfigure}%
\includegraphics[scale=1]{%
gesamttex/edit_VIII-4/images/LH_35_10_13_001-002_d6_002r.pdf}
\edcaption{%
\lbrack\textit{Fig.~6}\rbrack}
\end{edfigure}
\pstart
Explicandum superesset, quid fiat si
\edtext{pluribus quam duabus regulis}{%
\lemma{pluribus}%
\Bfootnote{%
\textit{(1)}~planis
\textit{(2)}~quam duabus regulis~\textit{L}}}
simul incumbat globus.\indexS{globus pluribus quam duabus regulis incumbens} 
Tunc autem ponendus est premere unam 
\edtext{regulam, perinde}{%
\lemma{regulam,}%
\Bfootnote{%
\textit{(1)}~quasi
\textit{(2)}~perinde~\textit{L}}}
ac si 
\edtext{globus\indexS{globus}}{%
\lemma{}%
\Bfootnote{%
globus 
\textit{erg.~L}}}
a 
\edtext{caeteris ipsa remota sustentaretur. Haec}{%
\lemma{caeteris}%
\Bfootnote{%
\textit{(1)}~sus
\textit{(2)}~sine ipsa sust
\textit{(3)}~ipsa remota sustentaretur
\textit{(a)}~, quo casu definiri potest, in qua recta descendere
\textit{(aa)}~conatur
\textit{(bb)}~conaretur
\textit{(b)}~. Haec recta 
\textit{(c)}~. Respe
\textit{(d)}~, et respectu \bars{ejus quae \textit{erg.}} pressionis regulae
\textit{(e)}~qu
\textit{(f)}~perind
\textit{(g)}~. Haec pressio sit e
\textit{(h)}~. Haec~\textit{L}}}
regula sit illa cujus pressio quaeritur. Porro videndum est in qua 
\edtext{recta conaturus}{%
\lemma{}%
\Bfootnote{%
recta \bars{et o \textit{erg. u. gestr.}} conaturus~\textit{L}}}
esset descendere globus si ipsa 
\edtext{abesset\lbrack,\rbrack\ in linea}{%
\lemma{abesset}%
\Bfootnote{%
\textit{(1)}~. Haec recta 
\textit{(a)}~sumat
\textit{(b)}~(quae cadet in idem planum cum
\textit{(c)}~deb
\textit{(d)}~\bars{erit ea quae \textit{streicht Hrsg.}}
\textit{(aa)}~in
\textit{(bb)}~glob
\textit{(cc)}~in plano
\textit{(dd)}~\bars{est \textit{erg.}} in plano communi globo et regulae cujus pressio
\bars{quaeritur \textit{streicht Hrsg.}}
globum tangit. Et perinde premetur haec regula, ac si recta illa esset regula nova globum cum ipsa consustentans
\textit{(2)}~in
\textit{(a)}~ea
\textit{(b)}~linea~\textit{L}}}
scilicet compositi conatus\indexS{conatus compositus} ex conatibus regularum 
\edtext{caeterarum.}{%
\lemma{}%
\Bfootnote{%
caeterarum 
\textit{erg.~L}}}
Haec linea si in planum 
\edtext{regulae propositae,}{%
\lemma{regulae}%
\Bfootnote{%
\textit{(1)}~et centri
\textit{(2)}~propositae,~\textit{L}}}
et centri,\indexS{planum per regulam et centrum} non incidit, fingatur globus 
\edtext{ab ea linea}{%
\lemma{ab ea}%
\Bfootnote{%
\textit{(1)}~per cent
\textit{(2)}~linea~\textit{L}}}
tanquam rigida perforatus, et secundum 
\edtext{eam descendere}{%
\lemma{eam}%
\Bfootnote{%
\textit{(1)}~descendens
\textit{(2)}~descendere~\textit{L}}}
tendens, qua ratione premat regulam. Is conatus cum non sit centralis,\indexS{conatus non centralis}
\edtext{seu non sit in plano per globum et regulam\lbrack,\rbrack\indexS{planum per regulam et globum}}{%
\lemma{}%
\Bfootnote{%
seu non sit in 
\textit{(1)}~glo
\textit{(2)}~plano per globum et regulam \textit{erg.~L}}}
compositus\indexS{conatus compositus} rursus intelligi potest,
\edtext{ex conatu tum}{%
\lemma{ex}%
\Bfootnote{%
\textit{(1)}~parallelo
\textit{(2)}~conatu
\textit{(a)}~ex centro
\textit{(b)}~tum~\textit{L}}}
\edtext{in plano secante normaliter in regula ipsa, planum}{%
\lemma{in plano}%
\Bfootnote{%
\textit{(1)}~per
\textit{(2)}~verticali
\textit{(3)}~normali ad planum in
\textit{(a)}~puncto contactus per
\textit{(b)}~regula ipsa,
\textit{(4)}~secante
\textit{(5)}~secante normaliter in regula ipsa,
\textit{(a)}~globum
\textit{(b)}~planum~\textit{L}}}
per regulam et centrum globi,\indexS{planum per regulam et centrum} tum in plano ipso globi, 
ita ut ex his duobus conatibus oriatur compositus\indexS{conatus compositus} in quo globus tendit.
\edtext{Et conatus qui}{%
\lemma{Et}%
\Bfootnote{%
\textit{(1)}~pre
\textit{(2)}~pressio quae
\textit{(3)}~conatus qui~\textit{L}}}
est in plano regulae cum globo, solum pro premente habendus.
\edtext{Tanquam linea}{%
\lemma{Tanquam}%
\Bfootnote{%
\textit{(1)}~is esset
\textit{(2)}~linea~\textit{L}}}
ejus esset altera regula, in qua globus praeter propositam incumberet. 
Reduco ergo conatum 
\edtext{globi pluribus}{%
\lemma{globi}%
\Bfootnote{%
\textit{(1)}~in
\textit{(2)}~pluribus~\textit{L}}}
planis simul 
\edtext{incumbentis,\indexS{conatus globi pluribus planis incumbentis} 
ita tamen ut ab ipsis solis non retineatur\lbrack,\rbrack}{%
\lemma{incumbentis}%
\Bfootnote{%
\textit{(1)}~(ita \lbrack...\rbrack\  retineatur)
\textit{(2)}~,~ita \lbrack...\rbrack\ retineatur;~\textit{L}}}
ad conatum in linea rigida per centrum 
\edtext{perforante; et hunc}{%
\lemma{perforante}%
\Bfootnote{%
\textit{(1)}~ad eum conatum
\textit{(2)}~, porro reduco ad eum
\textit{(3)}~; et hunc~\textit{L}}}
conatum cum 
\edtext{pos\lbrack s\rbrack it}{%
\lemma{}%
\Bfootnote{%
posit 
\textit{L ändert Hrsg.}}}
esse in plano diverso ab eo plano quod per regulam et globum transit\indexS{planum per regulam et globum}
reduco ad conatum componentem\indexS{conatus componens} 
\edtext{sumendum}{%
\lemma{}%
\Bfootnote{%
sumendum 
\textit{erg.~L}}}
in ipso plano per regulam 
\edtext{et globum.\indexS{planum per regulam et globum} Et fieri potest ut}{%
\lemma{et globum.}%
\Bfootnote{%
\textit{(1)}~Et
\textit{(2)}~Si plures glo
\textit{(3)}~Si plura plana \bars{sint \textit{streicht Hrsg.}} talia ut uno sublato non
\textit{(4)}~Et fieri potest 
\textit{(a)}~ut 
\textit{(aa)}~ipsa regula
\textit{(bb)}~ipsum regula
\textit{(b)}~ut~\textit{L}}}
nullius regulae planum per ipsam et centrum globi\indexS{planum per regulam et centrum globi} sit verticale.
Fieri etiam potest 
\edtext{ut una aut etiam plures regulae sublatae}{%
\lemma{ut}%
\Bfootnote{%
\textit{(1)}~plura aliquot
\textit{(2)}~unum aut etiam plura plana sublata
\textit{(3)}~una \lbrack...\rbrack\ sublatae~\textit{L}}}
non dent globo descendendi 
\edtext{potestatem. Tunc}{%
\lemma{potestatem.}%
\Bfootnote{%
\textit{(1)}~Itaque plura simul pr
\textit{(2)}~Videndum an combinationes institui possunt, qua
\textit{(3)}~Tunc duae re
\textit{(4)}~Tunc~\textit{L}}}
plures illae regulae sunt conjungendae et pro reliquis assumenda recta quam componunt, quasi perforans
\edtext{globum, haec}{%
\lemma{globum,}%
\Bfootnote{%
\textit{(1)}~et tunc 
\textit{(2)}~opino
\textit{(3)}~hae deinde re
\textit{(4)}~haec~\textit{L}}}
recta tanquam nova regula
\edtext{cum iis}{%
\lemma{cum}%
\Bfootnote{%
\textit{(1)}~caeteris
\textit{(2)}~iis~\textit{L}}}
quas simul tollere oportuit, est conjungenda.
\pend
\clearpage
\pstart\noindent
\href{\Scan{LH_35_10_13_002v}}%
{\lbrack2~v\textsuperscript{o}\rbrack}
\pend
\nobreak\begin{edfigure}[0.5em]%
\includegraphics[scale=1]{%
gesamttex/edit_VIII-4/images/LH_35_10_13_001-002_d7_002v.pdf}
\edcaption{%
\edtext{\lbrack\textit{Fig.~7}\rbrack}{%
\lemma{\lbrack\textit{Fig.~7}\rbrack}%
\Cfootnote[nonum]{%
Zwei gestrichene Entwürfe zum Diagramm werden nicht wiedergegeben.}}}
\end{edfigure}
\pstart
Si globus\indexS{globus in regula inclinata decurrens} \textit{A} decurreret in regula 
\edtext{inclinata}{%
\lemma{}%
\Bfootnote{%
inclinata 
\textit{erg.~L}}}
\edtext{\textit{BC}, quae tamen cedere}{%
\lemma{\textit{BC},}%
\Bfootnote{%
\textit{(1)}~quaeritur
\textit{(2)}~quae \bars{tamen \textit{erg.}}
cedere~\textit{L}}}
possit prementi, in angulo ad ipsam recto, quaeritur quantum oporteat esse 
\edtext{pondus \textit{I}, ut}{%
\lemma{pondus \textit{I},}%
\Bfootnote{%
\textit{(1)}~ut cessio
\textit{(2)}~ut~\textit{L}}}
cedere prohibeatur. Quoniam durante cessione globus \textit{A} semper manet in 
\edtext{\textit{BC}. Esto}{%
\lemma{\textit{BC}}%
\Bfootnote{%
\textit{(1)}~, patet
\textit{(2)}~, cedeat
\textit{(3)}~. Pon
\textit{(4)}~. Esto~\textit{L}}}
ergo dum \textit{BC} pervenit in (\textit{B})(\textit{C}) etiam \textit{A} pervenisse 
\edtext{in (\textit{A}). Etsi recta}{%
\lemma{in (\textit{A}).}%
\Bfootnote{%
\textit{(1)}~Q
\textit{(2)}~Ponendo ergo
\textit{(3)}~Etsi
\textit{(a)}~linea fortasse
\textit{(b)}~recta~\textit{L}}}
 \textit{A}(\textit{A}) non ideo asseratur esse verticalis,
cum globus\indexS{globus} simul duos motus, unum 
\edtext{deprimendi regulam,}{%
\lemma{deprimendi}%
\Bfootnote{%
\textit{(1)}~globum,
\textit{(2)}~regulam,~\textit{L}}}
alterum decurrendi in regula exercere possit. 
Dico igitur si quidem nullus sit globus\indexS{globus} \textit{I}, nec \textit{K},
globum \textit{A} descendere in recta verticali. Sed si absit globus \textit{K}, 
adsit globus \textit{I}, sed non satis 
\edtext{resistens nec}{%
\lemma{resistens}%
\Bfootnote{%
\textit{(1)}~ipsi \textit{A}
\textit{(2)}~nec~\textit{L}}}
sufficiens ad impedendam cessionem, centrum gravitatis\indexS{centrum gravitatis} amborum
\edtext{\textit{G}}{%
\lemma{}%
\Bfootnote{%
\textit{G} 
\textit{erg.~L}}}
recta verticali 
\edtext{\lbrack descendere\rbrack}{%
\lemma{}%
\Bfootnote{%
descendere
\textit{erg.~Hrsg.}}}
debere,
\edtext{seu \textit{G}(\textit{G})}{%
\lemma{seu}%
\Bfootnote{%
\textit{(1)}~\textit{\ind{1}G\ind{2}G}
\textit{(2)}~\textit{G}(\textit{G})~\textit{L}}}
esse 
\edtext{verticalem et}{%
\lemma{verticalem}%
\Bfootnote{%
\textit{(1)}~. Et quidem
\textit{(2)}~. Nempe recedente
\textit{(3)}~et~\textit{L}}}
quidem possibilium maximam. Sit 
\edtext{enim datus}{%
\lemma{enim}%
\Bfootnote{%
\textit{(1)}~ascensus
\textit{(2)}~datus~\textit{L}}}
ascensus ex \textit{I} 
\edtext{in (\textit{I}) assumtus}{%
\lemma{in (\textit{I})}%
\Bfootnote{%
\textit{(1)}~datus adeoque et
\textit{(2)}~assumtus~\textit{L}}}
eique aequalis $\beta(\beta)$ recessus regulae, sumatur \textit{\ind{2}G} pro arbitrio in 
\edtext{verticali \textit{G}(\textit{G}), et juncta (\textit{I})(\textit{G})}{%
\lemma{verticali}%
\Bfootnote{%
\textit{(1)}~\textit{\ind{1}G\ind{2}G}
\textit{(2)}~\textit{G}(\textit{G}),
et juncta
\textit{(a)}~\textit{\ind{2}G}
\textit{(b)}~(\textit{I})\textit{\ind{2}}
\textit{(c)}~(\textit{I})(\textit{G})~\textit{L}}}
occurret ipsi (\textit{B})(\textit{C}) in (\textit{A}).
\edtext{Oportet esse}{%
\lemma{Oportet}%
\Bfootnote{%
\textit{(1)}~\textit{G} sic assumi
\textit{(2)}~esse~\textit{L}}}
\textit{AG} ad \textit{GI}, ut (\textit{A})(\textit{G}) ad 
\edtext{(\textit{G})\lbrack(\textit{I})\rbrack}{%
\lemma{}%
\Bfootnote{%
(\textit{G})\textit{I}
\textit{L ändert Hrsg.}}}
unde dato pondere \textit{I} datur punctum (\textit{G}).
\pend
\pstart
\edtext{%
\edtext{Adda\lbrack ntur\rbrack}{%
\lemma{}%
\Bfootnote{%
Adda \textit{L ändert Hrsg.}}}%
}{%
\lemma{Adda\lbrack ntur\rbrack}%
\Cfootnote{%
Wohl abbrechender Anlauf zu der, anschließend auf der ersten Seite des Bogens nachgetragenen,
Anmerkung \glqq Addantur quae Bernoullius...\grqq, hier \refnote{FnA:35_10_13_001-002_Bernoullius}.}}
\lbrack\textit{Text bricht ab.}\rbrack
\pend
